\documentclass[a4paper,twoside]{article}

\usepackage{morewrites}

\usepackage[svgnames,dvipsnames,x11names]{xcolor}
\usepackage{graphicx}
\usepackage{texsmith-lists}
\usepackage[english]{babel}
\usepackage[colorlinks=true, linkcolor=NavyBlue, urlcolor=NavyBlue, citecolor=NavyBlue]{hyperref}
\usepackage[a4paper]{geometry}
\usepackage{ts-fonts}
\usepackage{csquotes}
\usepackage{amsmath}
\usepackage{float}
\usepackage{seqsplit}
\usepackage{ts-callouts}
\usepackage{ts-code}

\usepackage{lastpage}
\usepackage{refcount}
\makeatletter
\def\ps@plain{
\let\@oddhead\@empty
\let\@evenhead\@empty
\def\@oddfoot{
\hfil
\ifnum\getpagerefnumber{LastPage}>1\relax
\thepage
\fi
\hfil}
\let\@evenfoot\@oddfoot
}
\makeatother

\title{Headings}
\author{}\date{}
\begin{document}
\maketitle
\tscodeinline{\#} to \tscodeinline{\#\#\#\#\#\#}, six levels, \textbf{never numbered by hand}. Numbering and the
mapping to \LaTeX{} sectioning come from \tscodeinline{press.base\_level} and the template.
Attributes go at the end of the line.
\begin{tscode}[lang=md, engine=pygments]
\begin{Verbatim}[commandchars=\\\{\},breaklines, breakanywhere, commandchars=\\\{\}]
\PY{g+gh}{\PYZsh{} Firmware review}

\PY{g+gu}{\PYZsh{}\PYZsh{} Boot sequence \PYZob{}\PYZsh{}sec:boot\PYZcb{}}
\end{Verbatim}
\end{tscode}
Headings are \emph{relative}: TeXSmith realigns a multi-file hierarchy on its own ---
a per-fragment offset from the shallowest heading, plus the template slot base,
plus \tscodeinline{press.base\_level} --- and promotes the first heading to the document title
unless \tscodeinline{title:} is declared. That is processing, not syntax.
\section{Implicit ids}
A heading without \tscodeinline{\{\#id\}} still has an id, so the link forms resolve to it as
they do on GitHub and on every Python-Markdown site:
\begin{tscode}[lang=md, engine=pygments]
\begin{Verbatim}[commandchars=\\\{\},breaklines, breakanywhere, commandchars=\\\{\}]
\PY{g+gu}{\PYZsh{}\PYZsh{} Boot sequence}

See [](\PYZsh{}boot\PYZhy{}sequence) and [\PY{n+nt}{the boot sequence}](\PY{n+na}{\PYZsh{}boot\PYZhy{}sequence}).
\end{Verbatim}
\end{tscode}
The rule is GitHub's: take the plain text of the heading (inline markup and
code reduced to their text, zero-width nodes and the attribute list removed),
NFC-normalise it, lower-case it, drop every character that is not a letter, a
digit, a combining mark, a space, \tscodeinline{-\allowbreak{}} or \tscodeinline{\_}, and turn runs of spaces into one
\tscodeinline{-\allowbreak{}}. A duplicate gets \tscodeinline{-\allowbreak{}1}, \tscodeinline{-\allowbreak{}2}, ... in document order. Accents and non-Latin
scripts survive, as they do on GitHub and MkDocs Material.

The implicit id is a label like any other --- \tscodeinline{@boot-\allowbreak{}sequence} renders
\enquote{section 2} --- but it is derived at resolution, never stored and never printed,
so the round-trip is untouched. An explicit \tscodeinline{\{\#id\}} replaces it: a heading has
one id.
\begin{tscallout}[kind=warning, title={A reference to an implicit id is hinted}]
Editing the title changes the id and silently breaks every reference to it,
so referring to one raises the hint \tscodeinline{ref-\allowbreak{}implicit-\allowbreak{}id}, which suggests you
write \tscodeinline{\{\#id\}}. A site whose slugifier is not GitHub's (Python-Markdown's
default \tscodeinline{toc} drops accents) needs an explicit id anyway --- which is the
recommendation in any case.
\end{tscallout}
\section{Unnumbered and unlisted}
Two classes, both Pandoc's, each applying to its own heading only:
\begin{tscode}[lang=md, engine=pygments]
\begin{Verbatim}[commandchars=\\\{\},breaklines, breakanywhere, commandchars=\\\{\}]
\PY{g+gh}{\PYZsh{} Preface \PYZob{}.unnumbered\PYZcb{}}

\PY{g+gu}{\PYZsh{}\PYZsh{} Colophon \PYZob{}.unlisted\PYZcb{}}
\end{Verbatim}
\end{tscode}
\begin{description}
\item[{\tscodeinline{.unnumbered}}] Takes the heading out of the numbering sequence --- \tscodeinline{\textbackslash{}section*},
\tscodeinline{numbering: none}, \tscodeinline{class="unnumbered"}.
\item[{\tscodeinline{.unlisted}}] Does that \emph{and} keeps the heading out of the table of contents.
\end{description}
The document-level default is \tscodeinline{press.numbered}, true by default, which either
class overrides one heading at a time.
\begin{tscallout}[kind=note, title={\tscodeinline{\{-\allowbreak{}\}} is not an attribute list}]
Pandoc's shorthand for an unnumbered heading relies on bare-word
attributes, which TMark does not have (see Attributes), so
\tscodeinline{\{-\allowbreak{}\}} stays literal text. The Pandoc importer rewrites it.
\end{tscallout}
\section{Anchors and references}
A heading's \tscodeinline{\{\#id\}} is an anchor like any other, and the \textbf{host
decides the counter}: a heading is a section, so \tscodeinline{@sec:boot} renders
\enquote{section 2}. A prefix is required only where no host tells the kind --- or where
you want the heading numbered in a custom series instead of its
own:
\begin{tscode}[lang=md, engine=pygments]
\begin{Verbatim}[commandchars=\\\{\},breaklines, breakanywhere, commandchars=\\\{\}]
\PYZhy{}\PYZhy{}\PYZhy{}
press:
  declare:
    counters:
\PY{g+gu}{      fw: \PYZob{}name: Finding, format: \PYZdq{}FW\PYZhy{}\PYZob{}n:02d\PYZcb{}\PYZdq{}\PYZcb{}}
\PY{g+gu}{\PYZhy{}\PYZhy{}\PYZhy{}}

\PY{g+gu}{\PYZsh{}\PYZsh{} Boot loop \PYZob{}\PYZsh{}fw:boot\PYZhy{}loop\PYZcb{}}

The watchdog issue (@fw:boot\PYZhy{}loop) is fixed in 1.4.2.
\end{Verbatim}
\end{tscode}
See References for the reference forms.
\end{document}
